\section{Related Work}

This paper sits between four bodies of prior work.

First are studies and source documents that document inconsistency in vulnerability data sources. VIEM \cite{VIEM2019} shows that extracting affected software and versions from unstructured reports is hard and that agreement with NVD is far from complete. Croft et al. \cite{Croft2022} show that severity differs across sources and that those differences can affect downstream prediction. Croft et al. are the closest structured severity study, but their scope is one project/source family and they do not type non-conflict differences or adjudicate source support. NVD-focused audits such as "Cleaning the NVD" \cite{Anwar2022} and Zhang et al.'s \emph{The Flaw Within} \cite{Zhang2023} show that even a single database is not immune to quality problems, while user studies of NVD usage and CVSS scoring inconsistency show that practitioners encounter missing or incorrect entries and that severity scoring can be subjective even among evaluators \cite{NVDUsers2024,Wunder2024}. NVD process documents and recent operations updates further clarify that enrichment is an active analysis process, not a mechanical copy of CVE records, and that enrichment priorities can change under record-volume pressure \cite{NVDProcess2026,NVDStatus2026,NVDTransition2024,NVDOperations2026}. Recent affected-version benchmarking further shows that identifying affected versions is difficult even when tools use commits and advisories \cite{Chen2025AffectedVersions}. These works motivate the problem, but they do not directly solve post-alignment discrepancy typing for structured NVD-GHSA field pairs or evidence-constrained adjudication of residual factual conflicts. Unlike schema-cleaning or single-database repair work, this paper does not overwrite records or produce a canonical cleaned vulnerability database; it preserves both source values and emits a discrepancy type plus routing action.

Second are papers that detect inconsistency in unstructured vulnerability text. Sun et al.'s TOSEM work on aspect-level discrepancies \cite{Sun2023}, Wang et al.'s aspect extraction work on heterogeneous threat intelligence \cite{Wang2025Aspects}, and Li et al.'s 2025 \emph{VuldiffFinder} \cite{Li2025VuldiffFinder} are the closest neighbors in spirit, but they operate on report text and aspect extraction, not on already aligned structured NVD-GHSA fields. GapFinder \cite{GapFinder2021} is similar in method family but is aimed at broader security-intelligence text rather than aligned vulnerability records. VuldiffFinder is especially important because it is a recent Computers \& Security paper on vulnerability inconsistency; the boundary is that it discovers inconsistencies in unstructured vulnerability information, while this paper starts after database normalization/CVE alignment and studies type-first structured-field discrepancy handling. In this setting, the method output is not just an inconsistency flag: structured fields introduce schema-specific non-conflict mismatch modes such as subset reference lists, same-day timestamp representations, and range-bound shape differences, so the workflow emits an operational discrepancy type and routing action before any source-support question is asked.

Third are generic truth-discovery and multi-source fusion frameworks. CRH \cite{CRH2016} and truth-discovery surveys \cite{Wang2024TruthSurvey} provide methodological background for resolving conflicts across heterogeneous sources. Standard truth-discovery frames conflict resolution as true-value/source-reliability estimation; this setting instead uses explicit advisory evidence, field semantics, and an abstain/manual-review outcome. A CRH-style global reliability model is not the right first operation here because the main comparison has two source families and many post-alignment mismatches are representation, incompleteness, or timing differences rather than competing truth claims. Selective classification work \cite{Goren2024HSC} is relevant because it formalizes abstention under uncertainty, although this paper uses evidence and vulnerability-field semantics rather than learned classifier confidence.

Finally, vulnerability-intelligence ecosystem studies explain why NVD-GHSA discrepancies matter operationally. GHSA documentation and repository metadata show that the advisory database combines GitHub-reviewed advisories, unreviewed submissions, NVD imports, OSV-format records, and GitHub-specific affected-version extensions \cite{GHSAAbout2026,GHSARepo2026}. Work on the GHSA review pipeline highlights timing differences in advisory processing \cite{Segal2026GHSA}, VulZoo illustrates the scale of multi-source vulnerability intelligence integration \cite{Ruan2024VulZoo}, and scanner/evaluation studies show that heterogeneous sources, identifier schemes, and ambiguous version ranges can affect downstream vulnerability detection results \cite{Churakova2026VEX,Mandl2026VulnDetection}. These works do not solve field-level adjudication, but they motivate temporal discrepancies, source coverage gaps, affected-version ambiguity, and downstream scanner divergence as practical concerns.

The table below summarizes the boundary in reviewer-facing terms. The point is not that structured fields are easier than text; rather, structured fields fail in different ways after alignment, so a method needs field semantics, discrepancy types, and an abstention-capable evidence handoff.

\begin{table}[!htbp]
\centering
\setlength{\tabcolsep}{3pt}
\scriptsize
\caption{Reviewer-facing boundary between this paper and neighboring lines of work. The current contribution is a post-alignment structured-field typing workflow with an evidence-constrained adjudication scaffold.}
\label{tab:08-related-table-1}
\begin{tabularx}{\textwidth}{>{\raggedright\arraybackslash}p{0.18\textwidth}>{\raggedright\arraybackslash}X>{\raggedright\arraybackslash}X>{\raggedright\arraybackslash}X>{\raggedright\arraybackslash}X}
\toprule
Line of work & Input and stage & Output usually studied & Evidence and abstention & Boundary relative to this paper \\
\midrule
Vulnerability database quality, schema cleaning, and inconsistency studies & Single-source audits, schema repair, or cross-source observations & Missingness, quality defects, canonicalized records, severity differences, affected-version difficulty & Usually descriptive or repair-oriented; source support is not the main output & Motivates the problem, but this paper keeps both source values and types post-alignment field discrepancies instead of producing a canonical clean database \\
Unstructured advisory or threat-intelligence discrepancy mining & Report text, advisory prose, extracted aspects & Textual or aspect-level inconsistency labels & Evidence is the input text; abstention is usually not a first-class outcome & Closest in spirit, but this paper starts after CVE alignment and outputs a structured-field type plus routing action before source-support adjudication \\
Generic truth discovery and data fusion & Multi-source claims before domain-specific field routing & True-value and source-reliability estimates & Evidence is often implicit in source agreement; abstention is not central & This paper does not learn a global truth model; it first removes non-conflict field mismatches, then uses fetched advisory evidence and abstention for residual conflicts \\
Affected-version and scanner-ecosystem work & Package identities, range events, commits, scanner outputs & Applicability errors, scanner divergence, version-range difficulty & Evidence may include commits and advisories, but not this field-typing workflow & Explains the hardest field; the current method remains diagnostic until package-aware and range-aware adjudication is validated \\
\bottomrule
\end{tabularx}
\end{table}

Our paper is narrower and more concrete: after CVE-ID alignment, it types field-level discrepancies on structured NVD-GHSA pairs, then prototypes evidence-constrained adjudication with abstention on sampled factual conflicts. The current results are not a replacement for full truth discovery and are not a claim that silver labels are gold. The defensible contribution boundary is this type-first, post-alignment workflow: structured-field discrepancy typing first, and evidence-constrained source-support analysis only for residual factual conflicts.
